\documentclass[10pt,a4paper]{amsart}
\usepackage[margin=19mm]{geometry}
\usepackage{amsmath,amssymb,mathtools}
\usepackage[scr=rsfs,cal=euler]{mathalfa}
\usepackage{xfrac}
\usepackage[unicode,hidelinks,bookmarks=false]{hyperref}
\newcommand{\ie}{i.e.\ }
\newcommand{\cf}{cf.\ }
\newcommand{\resp}{respectively\ }
\newcommand{\Subsection}{\subsection}
\providecommand{\nobreakdash}{\nobreak\hbox{-}}
\setlength{\emergencystretch}{2em}
\multlinegap0pt
\usepackage{bm}
\usepackage{bbm}
\usepackage{dsfont}
\usepackage{xspace}
\usepackage{cancel}
\usepackage{mathtools}
\usepackage{amsthm}
\makeatletter
\@ifundefined{theorem}{\newtheorem{theorem}{Theorem}}{}
\makeatother
\title[Partially Alternative Real Division Algebras With A Few Imag]{Partially Alternative Real Division Algebras With A Few Imaginary Units}
\author{Tianran Hua, Marina Tvalavadze}
\date{Source: arXiv:2507.04919v2 (CC BY 4.0)}
\begin{document}
\maketitle

\noindent\emph{Source and licence.} Partially Alternative Real Division Algebras With A Few Imaginary Units. Version arXiv:2507.04919v2 (CC BY 4.0), distributed under the Creative Commons Attribution 4.0 International licence, as recorded in the arXiv licence metadata for this exact version and shown on the abstract page https://arxiv.org/abs/2507.04919v2. Full rights notice: licenses/arxiv-2507.04919-CC-BY-4.0.txt.\par\smallskip

\noindent\textit{Editorial scope.} Counted unit: the Theorem whose source label is \texttt{isomcond} in the article (source0.tex lines 657--659, source label \texttt{isomcond}), delivered together with the article's own complete original proof of it (source0.tex lines 662--721, one \texttt{proof} environment of 60 lines). No mathematics is written, repaired or re-derived: statement and proof are the article's own text line for line. Required context retained uncounted: none. Every definition, display and label of those lines is retained unchanged. Omitted from this package and not redistributed: 1--656, 660--661, 722--895. Nothing inside a retained excerpt is omitted or altered: every retained body is delivered exactly as the article's lines are written, so no omission receipt of any kind accompanies this package. Citations: the retained excerpts carry no citation command at all, so no bibliography entry of the article is transcribed and no reference is redistributed. Reference adaptation: none was needed and none was made; every reference inside the delivered unit resolves inside it, so no unresolved reference or citation is produced. Printed slips kept literal: no printed word, symbol, hypothesis or number of the counted statement or of its proof was altered. Layout and class: the article's own class is \texttt{amsart} with packages including amsmath, amssymb, babel, xy, amsthm, enumerate, graphicx; the wrapper is the lane's pinned \texttt{amsart} editorial wrapper at 10pt on A4, which restates the theorem-like environments the retained text needs and adds the article's own macro definitions that the retained text needs (none), with extra wrapper packages (bm, bbm, dsfont, xspace, cancel, mathtools). The delivered numbering is the unit's own. Assets: no retained span contains a figure environment, a graphics-inclusion command, a PDF-page inclusion, a table or a TikZ picture; no image file is copied into this package, the package declares no asset dependency and no image file is redistributed with it. Licence: version arXiv:2507.04919v2 of this article is distributed under the Creative Commons Attribution 4.0 International licence, as recorded in the arXiv licence metadata for this exact version and shown on the abstract page https://arxiv.org/abs/2507.04919v2. A full rights notice is delivered at licenses/arxiv-2507.04919-CC-BY-4.0.txt. Characters: every retained excerpt is pure ASCII, so the delivered engine, XeLaTeX with its default Latin Modern text font, reports no missing character.
\begin{theorem} \label{isomcond}
Let $A_{g, h}$ and $A_{g', h'}$ be two partially alternative division algebras given by multiplication table $(T_d)$ corresponding to parameters $g, h$ and $g', h'$, respectively.  Then the algebras are isomorphic if and only if $g=g'$ and $h = \pm h'$.
\end{theorem}

\begin{proof}
Assume that  $A_{g, h}$ and $A_{g', h'}$ are isomorphic. Let $\{1, i, j, k\}$ and $\{1, i', j', k'\}$ be two bases of $A_{g, h}$ and $A_{g', h'}$ with respect to which the algebras 
have multiplication tables of type $(T_d)$. 

 Consider the following possibilities. If $h=0$, then by Lemma 3 $g=-1$ and $A_{g, h}$ is the real quaternion algebra. Hence, by isomorphism,
$A_{g', h'}$ is the real quaternion algebra too, and, therefore, $h'=0$, $g'=-1$. Thus, the required condition holds true. 

Consider the case when $h\neq 0$, and, in particular, $\mathcal A_{g, h}$ is not isomorphic to the real quaternion algebra. It follows that $h'\neq 0$ since, otherwise, ${\mathcal A}_{g', h'}$ would be isomorphic to $\mathbb H$.

We observe that $A_{g,h}$ is generated as an algebra by the basis elements $i, j$.

Let $f: A_{g, h}\to A_{g', h'}$ be an algebra isomorphism.  By the above remark, we have that $f({\mathcal I}_{{\mathcal A}_{g, h}}) \subseteq {\mathcal I}_{{\mathcal A}_{g', h'}}$. 
Recall that if $h'\neq 0$, then $ \mathcal I_{{\mathcal A}_{g', h'}} = \{  x i' + y j' \,|\, x^2 + y^2 = 1\}.$ 

This implies that  both $f(i)$ and $f(j)$ belong to $ \mathcal I_{{\mathcal A}_{g', h'}}$, and, hence, 
$f(i) = \beta' i' + \gamma' j'$ and $f(j) = \beta'' i' + \gamma'' j'$ where $\beta'^2 +\gamma'^2 = 1$ and $\beta''^2 +\gamma''^2 = 1$.  

Since $k=ji$, we have that 
\begin{align*}
&f(k) = f(ji) = f(j) f(i) = (\beta'' i' + \gamma'' j'') (\beta' i' + \gamma' j')\\
&= (-\beta'\beta'' - \gamma'\gamma'' )1 + (\beta'\gamma'' - \gamma'\beta'') k'.
\end{align*}

Recall that $k^2 = g1+hi$ in $\mathcal A_{g, h}$. Then applying $f$ to $g1 + hi$ and $k^2$ separately, we obtain: 
\begin{align*}
& f(g1+hi)= g 1 + h f(i) = g1 + h \beta' i' + h \gamma' j' \,\, \text{equals\,to}  \\
& f(k^2)= f(k)^2 =  ( (-\beta' \beta'' - \gamma'\gamma'') 1 + (\beta' \gamma'' - \gamma' \beta'') k' )^2.
\end{align*}

We note that the expression for $f(k^2)$ contains no terms involving the basis element $j'$. This immediately implies that  the term $h \gamma' j'$ in $f(g1+hi)$ must be zero too. Hence, $\gamma'=0$ as $h\neq 0$. Since 
$\beta'^2 +\gamma'^2 = 1$, we have that $\beta'^2=1$. 
Thus, 
\begin{align*}
&f(k^2)= f(k)^2 =  ( (-\beta' \beta'' - \gamma'\gamma'') 1 + (\beta' \gamma'' - \gamma' \beta'') k' )^2\\
&= (\beta')^2 (-\beta'' 1 + \gamma'' k')^2 = (\beta'')^2 - 2 \beta''\gamma'' k' + (\gamma'' )^2 k'^2\\
&=  (\beta'')^2 - 2 \beta''\gamma'' k' + (\gamma'' )^2 (g' 1 + h' i')\\
& = ( ( \beta'')^2 + (\gamma'')^2 g')1 + h' (\gamma'')^2 i'  - 2 \beta''\gamma'' k'.
\end{align*}

Comparing the latter with the expression for  $f(g1+hi)$, we obtain the following equations:

\begin{align*}
 & (\beta'')^2 + (\gamma'')^2 g' = g, \tag{1}\\
 & (\gamma'')^2 h' = h\beta', \tag{2} \\
& \beta'' \gamma'' = 0. \tag{3}
\end{align*}

If $\gamma''=0$, then as follows from the second equation, $\beta'=0$. This means that $f(i) = \beta' i + \gamma' j$ is zero. This is impossible.

Thus, $\beta''=0$. Immediately, $(\gamma'')^2 = 1$.  This yields the required conditions on the parameters: $h' = \pm h$ and $g' =g$. 

Besides, we showed that if $f: A_{g, h}\to A_{g', h'}$ is an algebra isomorphism, then $f(i) = \beta' i'$ and $f(j)= \gamma'' j'$ where $(\beta')^2 = (\gamma'')^2 = 1.$


Conversely, if  the conditions  $g=g'$ and $h=\pm h'$ hold true, then either $A_{g, h}$ and $A_{g', h'}$ have the identical multiplication tables or we can use the mapping $f:\mathcal A_{g, h}\to \mathcal A_{g', h'}$
defined by $f(i)=-i'$ and $f(j)=j'$ to map isomorphically one algebra onto another one.

The proof is complete.

\end{proof}

\end{document}
